\documentclass[10pt,a4paper]{amsart}
\usepackage[margin=19mm]{geometry}
\usepackage{amsmath,amssymb,mathtools}
\usepackage[scr=rsfs,cal=euler]{mathalfa}
\usepackage{xfrac}
\usepackage[unicode,hidelinks,bookmarks=false]{hyperref}
\newcommand{\ie}{i.e.\ }
\newcommand{\cf}{cf.\ }
\newcommand{\resp}{respectively\ }
\newcommand{\Subsection}{\subsection}
\providecommand{\nobreakdash}{\nobreak\hbox{-}}
\setlength{\emergencystretch}{2em}
\multlinegap0pt
\usepackage{amssymb}
\newtheorem{lemma}{Lemma}
\newtheorem{theorem}{Theorem}
\def\U{{\mathcal{U}}}
\title{All doubly stochastic solutions to $AXA = XAX$ when $A$ is a permutation matrix}
\author{Bogdan D. Djordjevi\'c, Neboj\v{s}a \v{C}. Din\v{c}i\'c, Yingshan Chen}
\date{Source: arXiv:2609.05577v1 (CC BY 4.0)}
\begin{document}
\maketitle

\noindent\emph{Source and licence.} All doubly stochastic solutions to $AXA = XAX$ when $A$ is a permutation matrix. Version arXiv:2609.05577v1 (CC BY 4.0), distributed by arXiv under the Creative Commons Attribution 4.0 International licence, http://creativecommons.org/licenses/by/4.0/, as recorded in the arXiv licence metadata for this exact version and shown on the abstract page https://arxiv.org/abs/2609.05577v1. Full licence notice: \texttt{licenses/arxiv-2609.05577-CC-BY-4.0.txt}.\par\smallskip

\noindent\textit{Editorial scope.} Counted unit: the article's theorem udmus (source0.tex lines 711--713). The article's complete original proof of it is delivered with it (source0.tex lines 715--766, one \texttt{proof} environment of 52 lines). No mathematics is written, repaired or re-derived: the statement and the proof are the article's own lines, in the article's own notation, and no retained line is substituted or re-flowed.  Retained uncounted as the source-local context the proof needs: lines 74--75; lines 596--605; lines 640--642; lines 684--686; lines 688--690; lines 693--699.  Nothing was omitted from any retained span, so the package carries no omission record.  No retained line contains a citation, so the wrapper carries no bibliography and no bibliography entry.  No retained line contains a figure environment, an image-inclusion command or drawing code, so no image is delivered and the package declares no asset dependency.  Editorial formatting only, in addition: an AMS \texttt{amsart} preamble that restates the article's own declarations and macros used by the retained lines, namely the environments \texttt{lemma}, \texttt{theorem} on separate counters, together with the standard packages those lines need.  The retained environments print in this document as Lemma 1, Theorem 1 in order of appearance, whereas the article prints the same items with the numbering of its own full document.  The complete native source of the article as distributed in the arXiv e-print for this exact version is bundled unchanged as \texttt{sources/209449/source0.tex}.
	For a given square matrix $A\neq0$ (usually called a coefficient matrix), the Yang-Baxter-like matrix equation (YBME) is given as \begin{equation}
		\label{YBME} AXA=XAX.\end{equation}

	\begin{equation}
		\label{diagA}
		A=\left[\begin{array}{cccccc}
			0 & 0 & 0 & \ldots & 0 & 1\\
			1 & 0 & 0 &\ldots & 0 & 0\\
			0 & 1 & 0 & \ldots & 0 & 0 \\
			\vdots & \vdots& \vdots& \vdots& \vdots &\vdots\\
			0 & 0 & 0 &\ldots & 1 & 0
		\end{array}\right].
	\end{equation}

		\begin{equation}
			\label{R} R_X=\max_{1\leq j\leq n} X(i,j)=\max_{1\leq j\leq n} X(j,i).
		\end{equation}

	\begin{equation}
		\label{Ij}I_{j}=\{i\in\{1,\ldots,n\}: X(i,j)=R_X\}.
	\end{equation}

	\begin{equation}
		\label{Ji}J_{i}=\{j\in\{1,\ldots,n\}: X(i,j)=R_X\},
	\end{equation}

	\begin{lemma}\label{emptycomplements} Let $X$ be a doubly stochastic solution to \eqref{YBME}, with the constant $R_X$ provided as \eqref{R}. Let the sets $I_j$ and $J_i$ be provided via \eqref{Ij} and \eqref{Ji}, respectively, for each $1\leq i,j\leq n$. The following statements are equivalent:
		\begin{itemize}
			\item[(a)] There exists an $i_0$ such that $J_{i_0}=\{1,\ldots,n\}$.
			\item[(b)] There exists a $j_0$ such that $I_{j_0}=\{1,\ldots,n\}$.
			\item[(c)] $X=U_n$.
		\end{itemize}
	\end{lemma}

	\begin{theorem}\label{udmus}
		Let $A$ be given as \eqref{diagA}, and let $X$ be a doubly stochastic solution to \eqref{YBME}. Then there exists an $n_0\in\{1,\ldots,n\}$ such that $X\in\U_n(n_0)$. Moreover, the value $R_X$ from \eqref{R} is exactly $R_X=\frac{1}{n_0}$.
	\end{theorem}

	\begin{proof}
		If $X$ is regular then it is a permutation matrix, and $n_0=R_X=1$.
		
		If $X=U_n$, then $n_0=n$, and by Lemma \ref{emptycomplements} $R_X=\frac{1}{n}$. 
		
		Assume that $X$ is a proper doubly stochastic solution to \eqref{YBME} such that $X\neq U_n$. Let $(i_0,j_0)$ be a position such that $X(i_0-1,j_0+1)=R_X$. Since $I_{j_0}^c$ and $J_{i_0}^c$ are nonempty, we have
		\begin{eqnarray}\label{RXIj0}
			\begin{aligned}
				R_X&=X(i_0-1,j_0+1)=\sum_{s=1}^nX(i_0,s+1)X(s,j_0)=\\
				&=\sum_{s\in I_{j_0}}X(i_0,s+1)R_X+\sum_{s\notin I_{j_0}}X(i_0,s+1)X(s,j_0)\Leftrightarrow\\
				&\left(1-\sum_{s\in I_{j_0}}X(i_0,s+1)\right)R_X=\sum_{s\notin I_{j_0}}X(i_0,s+1)X(s,j_0)\Leftrightarrow\\
				&\left(\sum_{s\notin I_{j_0}}X(i_0,s+1)\right)R_X=\sum_{s\notin I_{j_0}}X(i_0,s+1)X(s,j_0).
			\end{aligned}
		\end{eqnarray}
		By definition, each $X(s,j_0)$ is strictly smaller than $R_X$ whenever $s\in I_{j_0}^c$, therefore the latter equality is possible if and only if 
		\begin{equation}\label{Ij0}
			\left(\forall s\in I_{j_0}^c\right)\quad X(i_0,s+1)=0.
		\end{equation}
		At the same time, \eqref{RXIj0} reads 
		\begin{eqnarray}\label{RXJi0}
			\begin{aligned}
				R_X&=X(i_0-1,j_0+1)=\sum_{s=1}^nX(i_0,s+1)X(s,j_0)=\\
				&=\sum_{s+1\in J_{i_0}}R_X X(s,j_0)+\sum_{s+1\notin J_{i_0}}X(i_0,s+1)X(s,j_0)\Leftrightarrow\\
				&\left(1-\sum_{s+1\in J_{i_0}}X(s,j_0)\right)R_X=\sum_{s+1\notin J_{i_0}}X(i_0,s+1)X(s,j_0)\Leftrightarrow\\
				&\left(\sum_{s+1\notin J_{i_0}}X(s,j_0)\right)R_X=\sum_{s+1\notin J_{i_0}}X(i_0,s+1)X(s,j_0).
			\end{aligned}
		\end{eqnarray}
		By the same argument as before, $X(i_0,s+1)<R_X$, whenever $s+1\in J_{i_0}^c$,  therefore
		\begin{equation}\label{Ji0} \left(\forall s+1\in J_{i_0}^c\right)\quad  X(s,j_0)=0.
		\end{equation}
		By introducing the set
		\begin{equation}
			\label{Ji01}J_{i_0}^c-1=\{k-1: k\in J_{i_0}^c\}\equiv \{k+_n (n-1): k\in J_{i_0}^c\},
		\end{equation}
		where again the operations modulo $n$ take values in the set $\{1,\ldots,n\}$, it follows that
		$$s+1\in J_{i_0}^c\Leftrightarrow s\in J_{i_0}^c-1,$$
		and \eqref{Ji0} reads \begin{equation}
			\label{Ji0-1} \left(\forall s\in J_{i_0}^c-1\right)\quad  X(s,j_0)=0.
		\end{equation}
		By definition, whenever $s\in I_{j_0}^c$, we have $X(s,j_0)<R_X$, therefore from \eqref{Ji0-1} we conclude that
		\begin{equation}
			J_{i_0}^c-1\subset I_{j_0}^c.
		\end{equation}
		And from \eqref{Ij0} we get
		$$\left(\forall s\in J_{i_0}^c-1\right)\quad X(i_0,s+1)=0\Leftrightarrow \left(\forall s+1\in J_{i_0}^c\right)\quad X(i_0,s+1)=0.$$
		In other words, 
		\begin{equation}\label{Xi0j}
			X(i_0,j)=\begin{cases} R_X,\quad j\in J_{i_0},\\0,\quad j\in J_{i_0}^c.
			\end{cases}
		\end{equation}
		This concludes that $n_0$ exists and is equal to $n_0:=|J_{i_0}|$, and $R_X=\frac{1}{n_0}$. Since $R_X$ is uniform throughout the matrix $X$, and $i_0$ was an arbitrarily chosen row, the proof is complete.
	\end{proof}

\end{document}
